%% Analysis of the negative/partial results, the degradation-ratio pitfall, and limitations.
%% Included by sigconf.tex. Do not add a preamble here.

\section{Analysis and Limitations}
\label{sec:analysis}

\paragraph{What the negative and partial results delimit.}
Two mechanisms did not behave as their motivation suggested, and both boundaries are informative. Mining's
inertness under its natural default shows that "train on the worst case" is empty unless the search space
exceeds what random exposure already covers. Distillation's profile shows that a richer observation helps when
the student cannot fit the target, which is a statement about task noise rather than about observation-process
shift; the two goals are separable and should not be conflated because both happen to involve sparse views.
Conversely, augmentation's dominance is itself a finding: informing a model by \emph{showing} it varied
observation processes was consistently more reliable than constraining what its representation must look like.

\paragraph{When the ratio metric misleads.}
P12's cross-entropy \emph{improves} under mild shifts while its AUROC and AUPRC fall, so
$\Delta_{\mathrm{shift}}$ computed on cross-entropy is negative and meaningless there; we omit it in
Table~\ref{tab:main} and read P12 from the ranking metrics. More generally, because
Eq.~\ref{eq:degradation} normalises by a model's own clean score, a uniformly weaker model can post a better
ratio. On USHCN-1mo the baseline's ratio ($+1.9\%$) is close to the framework's ($+1.5\%$) while its shifted
error is higher everywhere, and its clean error exceeds the framework's error \emph{under shift}. Absolute
error remains the quantity of interest. \todo{recompute degradation on AUROC for binary tasks}

\paragraph{Limitations.}
Every number is single-seed; the protocol requires five, and the smaller differences in
Tables~\ref{tab:components} and \ref{tab:main} are within the range we would expect seeds to move. Observation
processes are synthetic re-observations of real records, which is what makes the targets well defined but is
weaker evidence than a shift between two real acquisition environments. Compositions are limited to pairs. The
two backbones are tied to tasks rather than crossed with them. USHCN is not min-max normalised by its published
pipeline, so its errors are in raw units and not comparable in magnitude to the other forecasting settings. The
observation-aware encoder of Section~\ref{sec:method} is implemented but not evaluated here.
